% English translation of checked_pages/esq_119_030.tex; source: cote 119, site 30, PDF page 31, typescript page 28.
% AI-translated by GPT-6 Astra (OpenAI, ChatGPT Pro mode).
\scanpage{119/30}
\mspage{28}

the isotopy classes (in a sense dictated unambiguously by the
context) of isomorphisms (or even of morphisms more general
than isomorphisms). This idea, taken up successfully
in Yves Ladegaillerie's thesis in particular (cf. the beginning of par. 3),
has served me systematically in all my subsequent reflections
on combinatorial topology, whenever it was a matter of formulating
precise theorems translating topological situations
into combinatorial ones. In the present situation,
my hope was to reach the formulation (and proof !)
of an equivalence theorem between two suitable isotopy categories,
one being the category of ``admissible pairs'' $(X,Y)$,
the other that of ``admissible triples'' $(U,Y,f)$, where $Y$ is a
manifold, $U$ a manifold with boundary, and $f : \partial U \longrightarrow Y$ a proper
smooth fibration. Moreover, of course, I hoped that such a statement,
modulo work of an essentially algebraic nature, would
extend of itself to a more sophisticated statement applying to
general stratified structures.

Very quickly it became apparent that there could be no question
of obtaining so ambitious a statement in the context of
topological spaces, because of the perpetual phenomena of
``wildness''. Already when $X$ itself is a manifold and $Y$ is
reduced to a point, one runs up against the difficulty that the cone over
a compact space $Z$ can be a manifold at its vertex without
$Z$ being homeomorphic to a sphere, or even being a manifold. It
was also clear that the contexts of more rigid structures
existing at the time, such as the ``piecewise linear'' context,
were likewise inadequate - one of the prohibitive reasons they
shared being that, for a pair $(U,S)$ consisting of a
``space'' $U$ and a closed subspace $S$, and a gluing
map $f : S \longrightarrow T$, they did not allow one to construct the corresponding
amalgamated sum. It was a few years later that I was informed
of Hironaka's theory of the sets he calls, I believe,
``semianalytic'' (real), which satisfy some of the essential
stability conditions (probably even all of them)

\margin{30}

